\section{Metodologia da Pesquisa}
\label{sec:3.1}

Quanto à sua natureza, este trabalho classifica-se como pesquisa aplicada, pois tem como finalidade produzir conhecimento voltado à solução de um problema prático e delimitado: a automação da alocação de grades horárias em Instituições de Ensino Superior. Segundo \citeonline{gil2002}, a pesquisa aplicada caracteriza-se por estar orientada à aplicação prática dos conhecimentos produzidos, buscando responder a problemas concretos de determinado contexto. Nesse sentido, o presente estudo não se restringe à discussão teórica sobre o problema de timetabling universitário ou sobre Algoritmos Genéticos, mas propõe, implementa e avalia soluções computacionais integradas ao sistema IGRIS.

O problema investigado envolve a construção automatizada de grades horárias acadêmicas a partir de um conjunto de entidades e restrições, como turmas, disciplinas, professores, disponibilidades, turnos, matrizes curriculares e vínculos entre docentes e componentes curriculares. Trata-se de um problema de natureza combinatória, no qual o número de possibilidades cresce rapidamente conforme aumenta a quantidade de elementos envolvidos. Por esse motivo, a aplicação de técnicas heurísticas e meta-heurísticas torna-se relevante, uma vez que essas abordagens buscam encontrar soluções adequadas em tempo computacional viável, ainda que não garantam necessariamente a solução ótima global.

Do ponto de vista dos objetivos, a pesquisa classifica-se como exploratória e descritiva. A dimensão exploratória está associada à investigação da viabilidade técnica de diferentes implementações de Algoritmos Genéticos aplicadas ao mesmo problema. Conforme \citeonline{gil2002}, pesquisas exploratórias têm como finalidade proporcionar maior familiaridade com o problema, tornando-o mais explícito e permitindo o aprimoramento de ideias. Neste trabalho, essa característica manifesta-se na análise de três abordagens distintas, denominadas Sunny, Morpheus e Rocket, desenvolvidas com decisões próprias de implementação, integração e execução.

Complementarmente, a pesquisa possui caráter descritivo, pois registra, organiza e analisa sistematicamente o comportamento das implementações desenvolvidas. De acordo com \citeonline{prodanov2013}, a pesquisa descritiva busca observar, registrar e analisar fenômenos sem manipulá-los de forma arbitrária. No contexto deste trabalho, essa classificação se justifica pela coleta estruturada de métricas de benchmark, como tempo de execução, consumo de recursos computacionais e qualidade das soluções geradas. Assim, a pesquisa não apenas apresenta os algoritmos implementados, mas descreve seu comportamento em um cenário experimental controlado.

Quanto à abordagem do problema, o estudo possui natureza quantitativa e qualitativa, com predominância quantitativa na etapa de avaliação experimental. A dimensão quantitativa está presente na coleta e comparação de métricas numéricas obtidas durante os benchmarks. Esses dados permitem avaliar o desempenho das implementações sob critérios mensuráveis, reduzindo a subjetividade da análise comparativa. Conforme \citeonline{prodanov2013}, a pesquisa quantitativa considera que os fenômenos podem ser traduzidos em números, opiniões e informações classificáveis e analisáveis.

Entretanto, o trabalho também apresenta uma dimensão qualitativa, especialmente nas etapas de análise arquitetural, modelagem do domínio acadêmico e escolha dos dilemas de implementação. A comparação entre Sunny, Morpheus e Rocket não se limita aos valores numéricos obtidos nos testes, pois também envolve a interpretação das decisões de projeto, da estrutura de integração, da organização dos módulos e das características técnicas de cada solução. Dessa forma, embora a etapa de benchmark forneça a base quantitativa da avaliação, a análise qualitativa contribui para compreender as razões técnicas que podem influenciar os resultados observados.

Em relação aos procedimentos técnicos, esta pesquisa combina revisão bibliográfica estruturada e pesquisa experimental. A revisão bibliográfica foi utilizada para fundamentar os principais conceitos abordados no trabalho, incluindo o problema de timetabling universitário, Algoritmos Genéticos, heurísticas, meta-heurísticas, arquitetura de software, dilemas de programação e trabalhos relacionados. Segundo \citeonline{lakatos2003}, a pesquisa bibliográfica permite colocar o pesquisador em contato com o conhecimento já produzido sobre determinado tema, oferecendo base teórica para a construção da investigação.

A etapa bibliográfica foi conduzida a partir da análise de livros, artigos científicos, publicações técnicas e estudos correlatos. Essa fundamentação permitiu delimitar o problema, identificar abordagens já utilizadas na literatura e justificar a escolha dos Algoritmos Genéticos como estratégia de otimização. Também contribuiu para embasar as decisões relacionadas à arquitetura do sistema, à separação dos componentes e à comparação entre diferentes formas de implementação.

A pesquisa experimental, por sua vez, materializa-se na implementação e avaliação das três soluções propostas. Conforme \citeonline{gil2002}, a pesquisa experimental envolve a determinação de um objeto de estudo, a seleção das variáveis que podem influenciá-lo e a definição de formas de controle e observação dos efeitos produzidos. Neste trabalho, o objeto de estudo corresponde ao desempenho das implementações Sunny, Morpheus e Rocket na geração automatizada de grades horárias. As variáveis observadas incluem métricas de desempenho computacional e qualidade da solução, enquanto as condições de execução foram padronizadas para permitir comparação entre os algoritmos.

Para reduzir interferências externas, os benchmarks foram conduzidos sob condições equivalentes de hardware, base de entrada e parâmetros de execução. Essa padronização buscou garantir que as diferenças observadas entre Sunny, Morpheus e Rocket estivessem associadas principalmente às características das implementações, e não a variações acidentais do ambiente de teste. Foram executadas múltiplas requisições para cada algoritmo, totalizando 9.000 requisições por implementação, conforme o planejamento experimental apresentado no Capítulo 7.

O desenvolvimento do sistema seguiu uma estratégia iterativa e incremental, inspirada em práticas do Kanban. Essa abordagem foi adotada para organizar a construção gradual dos componentes do IGRIS e das implementações dos algoritmos. Conforme \citeonline{anderson2010}, o Kanban propõe tornar o fluxo de trabalho visível, limitar o trabalho em progresso e promover melhorias evolutivas no processo. No contexto deste trabalho, essas práticas foram operacionalizadas por meio da divisão das atividades em etapas de acompanhamento, como tarefas planejadas, tarefas em desenvolvimento, tarefas em validação e tarefas concluídas. Também foi considerada a limitação explícita de trabalho em progresso, conhecida como WIP Limit, a fim de evitar a execução paralela excessiva de tarefas e favorecer a finalização incremental dos módulos. Com base nessa organização, o desenvolvimento ocorreu em ciclos sucessivos de planejamento, implementação, validação e ajuste, permitindo que os componentes fossem construídos e avaliados progressivamente antes da integração final.
